\documentclass[11pt]{article}
\usepackage[margin=1in]{geometry}
\usepackage[T1]{fontenc}
\usepackage{amsmath,amssymb,amsthm}
\usepackage{enumitem}
\usepackage{microtype}
\usepackage{booktabs}
\usepackage{xcolor}
\usepackage{hyperref}

\hypersetup{
  colorlinks=true,
  linkcolor=blue,
  citecolor=blue,
  urlcolor=blue,
  pdftitle={Dynamic Regular Witnesses for Near-Exact State Complexity of Explicit Binary Consensus with Periodic Time},
  pdfauthor={Ryutaro Yonezu}
}

\newtheorem{theorem}{Theorem}
\newtheorem{lemma}{Lemma}
\newtheorem{proposition}{Proposition}
\newtheorem{corollary}{Corollary}
\theoremstyle{remark}
\newtheorem{remark}{Remark}

\newcommand{\Fzero}{F_0}
\newcommand{\Fone}{F_1}
\newcommand{\Live}{L}

\title{\textbf{Dynamic Regular Witnesses for Near-Exact State Complexity\\
of Explicit Binary Consensus with Periodic Time}}
\author{Ryutaro Yonezu\\\small Independent Researcher}
\date{September 6, 2026\\\small Version v1.0.0 -- not peer reviewed}

\begin{document}
\maketitle

\begin{abstract}
We study deterministic binary consensus with explicit termination in anonymous synchronous
$1$-interval-connected dynamic networks under one-bit broadcast-counting communication.
Each round $t$ exposes a free globally aligned phase $\phi_t=t\bmod P$, while the full
round number is unavailable.  Let $S_n(n,P)$ denote the minimum number of persistent
states required when the network size $n$ and the phase period $P$ are known
non-uniformly.

A direct predecessor obtained the known-size bounds
\[
2+\left\lceil\frac{\lfloor n/2\rfloor}{P}\right\rceil
\le S_n(n,P)\le
2+\left\lceil\frac{n-1}{P}\right\rceil,
\]
leaving a factor-two gap in the numerator.  We reduce this to an additive constant.
The main idea is a \emph{dynamic regular witness}: every communication snapshot is
regular, so opposite homogeneous executions follow topology-independent phase-augmented
state orbits, while the topology is rearranged to make two backward causal cones grow
almost one vertex per round.  Using explicit quartic gadgets, for every $n\ge15$,
\[
\boxed{
2+\left\lceil\frac{n-13}{P}\right\rceil
\le S_n(n,P)\le
2+\left\lceil\frac{n-1}{P}\right\rceil.
}
\]
For even $n\ge12$, cubic witnesses improve the lower numerator from $n-13$ to $n-10$.
Thus, in the unsaturated regime where $n/P\to\infty$, the leading coefficient of the known-size state complexity is one.

We also determine the exact saturation threshold for the minimum possible three
persistent states:
\[
\boxed{
S_n(n,P)=3\quad\Longleftrightarrow\quad P\ge n-1
}
\qquad(n\ge4).
\]
The lower direction uses only a static path and combines anonymous path symmetry,
a forced maximal-speed zero wave, and the binary broadcast alphabet.
The claims are restricted to the stated model.  We do not claim the full exact formula
$S_n(n,P)=2+\lceil(n-1)/P\rceil$ for arbitrary state budgets.
\end{abstract}

\noindent\textbf{Keywords.}
anonymous dynamic networks; binary consensus; explicit termination; state complexity;
periodic time; one-bit communication; regular dynamic graphs; causal cones; local memory.

\section{Introduction}

Explicit termination forces an anonymous process to represent evidence that information has
had enough time to propagate.  This is different from stabilization: a binary consensus
value can be propagated by a two-state monotone flooding process, while a process that must
know when it is safe to stop needs additional local timing information.

The state--phase line of work studied this cost when every node receives the free common
phase
\[
\phi_t=t\bmod P
\]
but not the absolute round number.  In that model a direct predecessor proved, for known
network size $n$,
\[
2+\left\lceil\frac{\lfloor n/2\rfloor}{P}\right\rceil
\le S_n(n,P)\le
2+\left\lceil\frac{n-1}{P}\right\rceil,
\]
after improving an earlier split-cycle lower bound.  The upper bound is a block-counter
implementation of monotone zero flooding with horizon $n-1$.  The remaining question is
whether the factor-two gap in the lower numerator reflects the task or only the static-cycle
witness.

This paper shows that the latter is largely responsible.  Static cycles are convenient
because their homogeneous executions are symmetric, but their ordinary diameter is only
about $n/2$.  The key observation is that static symmetry is unnecessary.  If every
snapshot is $d$-regular, then a homogeneous execution is blind to the edge identities:
every node has the same state, broadcasts the same bit, and receives the same degree-$d$
count vector.  The adversary can therefore change the regular topology aggressively without
changing the homogeneous state orbit.

We exploit this freedom by constructing two disjoint backward causal cones.  Each cone has
a constant-size first expansion forced by regularity and then grows by only one vertex per
round.  Explicit path/cycle-plus-matching completions fill the rest of every snapshot while
preserving regularity and connectivity.  The resulting waiting-sum lower bound is
$n-O(1)$, which combines with two-orbit phase packing to give a near-exact state lower
bound with leading constant one.

A second result isolates the smallest state budget.  Three persistent states necessarily
have the normal form ``final zero, one live state, final one.''  On a static path this
collapses the pre-termination symmetry enough to determine exactly how much free phase is
needed: three states are possible precisely when the period is at least $n-1$.

\paragraph{Main contributions.}
\begin{enumerate}[leftmargin=1.5em]
\item A regular-homogeneous invariance lemma and a dynamic causal-cone construction that
decouple homogeneous state evolution from mixed-input information propagation.
\item For all $n\ge15$,
\[
2+\left\lceil\frac{n-13}{P}\right\rceil
\le S_n(n,P)\le
2+\left\lceil\frac{n-1}{P}\right\rceil.
\]
\item For even $n\ge12$, the stronger cubic lower bound
\[
S_n(n,P)\ge
2+\left\lceil\frac{n-10}{P}\right\rceil.
\]
\item The exact three-state saturation threshold
\[
S_n(n,P)=3\iff P\ge n-1
\qquad(n\ge4).
\]
\item An explicit finite construction audit over $6253$ complete parameter cases; the
computation is used only to test the constructive lemmas, not as a substitute for proof.
\end{enumerate}

\section{Relation to prior work}

The direct predecessor is Yonezu's two-orbit packing note~\cite{yonezu2026twoorbit},
which made the known-diameter state complexity exact and strengthened the known-size lower
bound to
$2+\lceil\lfloor n/2\rfloor/P\rceil$.  The present paper addresses the residual known-size
gap by replacing the static-cycle witness with regular dynamic witnesses and by isolating
the exact three-state threshold.

Parzych and Daymude study memory lower bounds for terminating broadcast in anonymous,
synchronous, $1$-interval-connected dynamic networks~\cite{parzych2026}.  Their results
show that termination memory is a genuine resource in nearby anonymous dynamic settings,
but the task, start semantics, and objective differ from the periodic-phase binary-consensus
state complexity considered here.

Blanc, Di Luna, and Viglietta study deterministic computation when every anonymous node
broadcasts one bit and observes only the numbers of neighboring zero and one broadcasts
\cite{blanc2026}.  Their results establish broad computational power and round-complexity
bounds under the same severe communication interface.  They do not give the persistent
state--phase thresholds studied here.

Di Luna and Viglietta develop finite-state and self-stabilizing computation in anonymous
dynamic networks~\cite{diluna2024}.  Their finite-state stabilization framework addresses
general computation rather than exact persistent-state requirements for explicitly
terminating binary consensus with an externally supplied phase.  Turau gives a randomized
$O(\log\log n)$-memory algorithm for anonymous dynamic broadcast with stabilizing
termination~\cite{turau2026}; randomization, the broadcast task, and stabilizing rather
than explicit termination place that result outside the present state-threshold question.

Modulo-clock synchronization is a different line.  Penet de Monterno, Charron-Bost, and
Merz synchronize clocks modulo $P$ in dynamic networks~\cite{monterno2023}, and
Charron-Bost and Penet de Monterno study self-stabilizing clock synchronization
\cite{charron2022}.  Those works compute or recover a clock signal; here the common
periodic phase is supplied externally for free, and the question is how many persistent
states explicit consensus still requires.

A targeted literature search through September 6, 2026 did not identify a prior theorem
giving either the near-exact $n-O(1)$ numerator bounds or the exact three-state threshold
for the stated model.  This is a scoped literature-screening statement, not a certification
of novelty.

\section{Model and definitions}

We use the same state--phase model as the direct predecessor on periodic-time explicit consensus. Fix $n$ and $P$.  Nodes are anonymous, deterministic, and execute identical code.
In propagation round $t$, the adversary chooses a simple undirected connected graph
$G_{t+1}$.  Each node broadcasts one bit to all current neighbors and observes only
the pair $(c_0,c_1)$ of received zero and one broadcasts.  The common free phase is
\[
\phi_t=t\bmod P.
\]
The broadcast and transition functions have the form
\[
\sigma:Q\times\{0,\ldots,P-1\}\to\{0,1\},
\]
\[
\delta:Q\times\{0,\ldots,P-1\}\times\mathbb N\times\mathbb N\to Q.
\]
The initialization is deterministic; write $\iota(b)\in Q$ for the persistent state
of an input-$b$ node at time $0$.  Finality and output depend only on persistent state.
Once a final state is entered, the persistent state and output never change, although
the node remains communication-present.

\section{Dynamic causal locality and regular homogeneous orbits}

For a target $v$ at time $T$, define its backward causal sets by
\[
B_T(v)=\{v\},
\qquad
B_t(v)=B_{t+1}(v)\cup N_{G_{t+1}}(B_{t+1}(v))
\quad(0\le t<T).
\]

\begin{lemma}[Dynamic causal-cone locality]
Fix a graph sequence and a target $v$ at time $T$.  If two executions have the same
initial persistent states on $B_0(v)$, then $v$ has the same state at every time
$0,\ldots,T$ in the two executions.
\end{lemma}

\begin{proof}
Induct backward through the dependency relation.  The state of a node at time $t+1$
depends only on its own state at time $t$, the common phase $\phi_t$, and the broadcast
bits of its neighbors in $G_{t+1}$.  Thus every dependency of the target at time $T$
is contained in the recursively defined backward cone.  The phase sequence is common
to the two executions, so equality of the initial states on $B_0(v)$ implies equality
of all states and broadcasts that can influence $v$ through time $T$.
\end{proof}

\begin{lemma}[Regular homogeneous invariance]
Fix $d\ge1$.  On every dynamic graph sequence whose snapshots are all $d$-regular,
a homogeneous input execution has one common persistent state $q_t$ at every node
at each time $t$.  Moreover, the augmented state $(q_t,\phi_t)$ evolves by a
deterministic map
\[
G_d:Q\times\{0,\ldots,P-1\}\to Q\times\{0,\ldots,P-1\}
\]
that is independent of the particular $d$-regular topology schedule.
\end{lemma}

\begin{proof}
If all nodes have the same state $q$ at phase $\phi$, they broadcast the same bit
$\sigma(q,\phi)$.  Every node has degree $d$, so every node receives the same count
vector, either $(d,0)$ or $(0,d)$.  Determinism therefore gives the same next state
at every node.  The observation depends only on $(q,\phi,d)$, not on the identities
of the incident edges or on the topology schedule.
\end{proof}

For $b\in\{0,1\}$, let $\tau_b$ be the first final time in the homogeneous input-$b$
execution on a $d$-regular schedule.  By the preceding lemma, $\tau_b$ depends on
the algorithm, $n,P$, and $d$, but not on the particular $d$-regular schedule.

\begin{lemma}[Regular two-orbit packing]
For the two opposite homogeneous executions on $d$-regular schedules,
\[
P|Q|\ge \tau_0+\tau_1+2P.
\]
\end{lemma}

\begin{proof}
The proof is the same phase-augmented packing argument used for homogeneous static
cycles.  Before the first final time, an augmented state cannot repeat within either
homogeneous orbit, or deterministic iteration of $G_d$ would loop and prevent first
termination.  The two pre-final augmented orbits are disjoint: if they met, the two
deterministic suffixes would coincide and eventually reach the same final output,
contradicting opposite unanimity requirements.  Finally, the distinct immutable
final states reached by the two homogeneous executions each sweep through all $P$
phase values.  Hence the augmented state space contains at least
$\tau_0+\tau_1+2P$ distinct pairs.
\end{proof}

\section{Explicit regular protected blocks}

\begin{lemma}[Quartic protected block]
Let $S$ be a vertex set of size $s\ge5$ and let $x\notin S$.  There is a simple
connected partial graph on $S\cup\{x\}$ such that every vertex of $S$ has degree $4$,
\[
N(S)\setminus S=\{x\},
\]
and $x$ has exactly two incident edges into $S$.
For $s=5$, any prescribed target $z\in S$ can be chosen to have all four neighbors
inside $S$.
\end{lemma}

\begin{proof}
Start from the square $C_s^2$, which is $4$-regular for $s\ge5$.
Delete one edge $uv$ and add $ux$ and $vx$.  All vertices of $S$ again have degree
$4$, while $x$ has two incident edges into $S$.  For $s=5$, $C_5^2=K_5$; choose
the deleted edge disjoint from the prescribed target $z$.
\end{proof}

\begin{lemma}[Cubic protected block]
Let $S$ be a vertex set of size $s\ge4$ and let $x\notin S$.
There is a simple connected partial graph in which every vertex of $S$ has degree $3$
and $N(S)\setminus S=\{x\}$.  If $s$ is even, $x$ has two incident edges into $S$;
if $s$ is odd, $x$ has one incident edge into $S$.
For $s=4$, any prescribed target $z\in S$ can be chosen to have all three neighbors
inside $S$.
\end{lemma}

\begin{proof}
If $s$ is even, take a cycle $C_s$ and add the opposite-vertex perfect matching,
pairing vertex $i$ with $i+s/2$ (indices modulo $s$).  These are non-cycle edges
for $s\ge4$, so the result is a connected cubic graph.  Delete one matching edge
$uv$ and add $ux,vx$.
For $s=4$ this is $K_4$ with one edge replaced by the two edges to $x$, and the
deleted edge can be chosen disjoint from the prescribed target.

If $s=2h+1$ is odd, take $C_s$ and a matching
\[
\{\, i(i+h): 1\le i\le h\,\}
\]
on the vertices $1,\ldots,2h$, leaving vertex $0$ unmatched.  The matching edges
are non-cycle edges because $h\ge2$.  Add the edge $0x$.  Every vertex of $S$
then has degree $3$, and $x$ has one edge into $S$.
\end{proof}

\section{Explicit completion of the complement}

\begin{lemma}[Quartic complement completion]
Let $R$ be a vertex set with $|R|\ge6$.  Mark at most two vertices of $R$ as
connectors.  Require each connector to have degree $2$ inside $R$ and every other
vertex to have degree $4$.  Then there is a simple connected graph on $R$ realizing
those degrees.
\end{lemma}

\begin{proof}
With no connector, use $C_{|R|}^2$.

With one connector $x$, write the remaining vertices as
$o_1,\ldots,o_q$, where $q\ge5$.  Use the cycle
\[
x,o_1,o_2,\ldots,o_q,x,
\]
so every vertex currently has degree $2$.  Add a Hamiltonian cycle on the ordinary
vertices that avoids the path edges $o_io_{i+1}$.  One explicit order is
\[
o_1,o_3,o_5,\ldots,o_2,o_4,o_6,\ldots,o_1,
\]
with the odd and even subsequences truncated in the natural way.  For $q\ge5$,
none of its edges is a path edge.  The ordinary vertices reach degree $4$ while
$x$ remains degree $2$.

With two connectors $x,y$ and $|R|\ge7$, use the base cycle
\[
x,o_1,\ldots,o_q,y,x
\]
and the same edge-disjoint Hamiltonian cycle on the $q\ge5$ ordinary vertices.
Both connectors have degree $2$ and all ordinary vertices degree $4$.

The only remaining case is $|R|=6$ with two connectors.  Start from $K_6$ and
delete the five edges
\[
xy,\ xa,\ xb,\ yc,\ yd,
\]
where $a,b,c,d$ are the four ordinary vertices.  The connectors then have degree
$2$ and all ordinary vertices degree $4$.  The resulting graph is connected.
\end{proof}

\begin{lemma}[Cubic complement completion]
Let $R$ be a vertex set with $|R|\ge5$.  Mark at most two vertices as connectors,
each with required degree $1$ or $2$ inside $R$, and require every ordinary vertex
to have degree $3$.  If the total required degree is even, then a simple connected
realization exists.
\end{lemma}

\begin{proof}
The parity condition is necessary by the handshake lemma.  We give explicit
realizations.

With no connector, the parity condition makes $|R|$ even.  Use a cycle plus a
perfect matching on opposite vertices.

With one connector $x$ of required degree $2$, use the cycle
\[
x,o_1,\ldots,o_q,x
\]
and add a perfect matching on the even number $q$ of ordinary vertices, using
non-path edges.  If $x$ has required degree $1$, use a cycle on the odd number
$q$ of ordinary vertices, attach $x$ to $o_1$, and add a perfect matching on
$o_2,\ldots,o_q$ using non-cycle edges.

With two connectors of required degrees $(1,1)$, use the path
\[
x,o_1,\ldots,o_q,y
\]
and add a perfect matching on the even number $q$ of ordinary vertices.
For degrees $(2,2)$, use the cycle
\[
x,o_1,\ldots,o_q,y,x
\]
and again add a perfect matching on the even number $q$ of ordinary vertices.
For degrees $(1,2)$ and $q\ge5$, use the cycle
\[
y,o_1,o_2,\ldots,o_q,y,
\]
attach $x$ to $o_1$, and add a perfect matching on
$o_2,\ldots,o_q$.  The parity condition makes $q$ odd.

In every matching step above, the vertices to be matched form an even ordered list
$w_1,\ldots,w_m$ with $m\ge4$, except for the special case below.  Pair
$w_i$ with $w_{i+m/2}$ for $1\le i\le m/2$.  Each pair is separated by at
least two positions in the displayed path or cycle, so no matching edge duplicates
an existing edge.  The only size-five case not covered by these formulas is
$q=3$ with connector degrees $(1,2)$.  On vertices $x,y,a,b,c$, take $K_5$ and
delete
\[
xy,\ xa,\ xb,\ yc.
\]
The remaining degrees are respectively $1,2,3,3,3$, and the graph is connected.
\end{proof}

\section{Slow disjoint causal cones}

\begin{proposition}[Quartic slow-cone schedule]
Let $n\ge15$ and let $t_0,t_1\ge0$ satisfy
\[
t_0+t_1\le n-14.
\]
There exists a sequence of connected $4$-regular graphs and two targets
$z_0,z_1$ such that their backward causal cones through horizons $t_0,t_1$
are disjoint at time $0$.
\end{proposition}

\begin{proof}
Choose disjoint vertex pools $U_b$ of size $t_b+4$, choose a target
$z_b\in U_b$ for each $b\in\{0,1\}$, and let
\[
M=V\setminus(U_0\cup U_1).
\]
The hypothesis gives $|M|\ge6$.

For each $b$ with $t_b>0$, construct nested sets
\[
C_b(t_b)=\{z_b\},
\]
\[
|C_b(t_b-1)|=5,
\]
and, for $0\le t\le t_b-2$,
\[
C_b(t)=C_b(t+1)\cup\{x_{b,t}\},
\]
using fresh vertices from $U_b$.  Thus $|C_b(0)|=t_b+4$.

We now choose the snapshot $G_{t+1}$ independently for each round $t$.
For each active target $b$ with $t<t_b$, impose one protected block.

If $t=t_b-1$, apply the quartic protected-block lemma to
$S=C_b(t)$.  Choose its external connector in $M$, distinct from the connector
of any other protected block in the same snapshot, and choose the deleted edge
disjoint from $z_b$.  Then all four neighbors of $z_b$ lie in $C_b(t)$, so the
one-step backward predecessor of $\{z_b\}$ is exactly $C_b(t)$.

If $t\le t_b-2$, apply the protected-block lemma to
$S=C_b(t+1)$ with external connector
$x_{b,t}=C_b(t)\setminus C_b(t+1)$.  Then
\[
C_b(t+1)\cup N_{G_{t+1}}(C_b(t+1))=C_b(t).
\]

The protected sets are disjoint, and their connectors are distinct and lie
outside both protected sets.  Remove the protected-set vertices from consideration
and let $R$ be the remaining vertices, including the at most two connectors.
Because each protected set lies inside its pool,
\[
|R|\ge |M|\ge6.
\]
Each connector already has two edges into its protected block and therefore needs
degree $2$ inside $R$; every other vertex of $R$ needs degree $4$.
Complete $R$ by the quartic complement-completion lemma.
The union is simple, connected, and $4$-regular.

For every positive horizon, induction on the backward-cone recursion gives that the
exact cone of $z_b$ at time $t$ is $C_b(t)$ for every relevant $t$.  When $t_b=0$
the time-zero cone is simply $\{z_b\}$.  Hence in all cases the two time-zero
cones are contained in the disjoint pools $U_0,U_1$.
\end{proof}

\begin{proposition}[Cubic slow-cone schedule]
Let $n\ge12$ be even and let $t_0,t_1\ge0$ satisfy
\[
t_0+t_1\le n-11.
\]
There exists a sequence of connected $3$-regular graphs and two targets whose
time-zero backward causal cones through horizons $t_0,t_1$ are disjoint.
\end{proposition}

\begin{proof}
Choose disjoint pools $U_b$ of size $t_b+3$, choose a target
$z_b\in U_b$ for each $b\in\{0,1\}$, and let
$M=V\setminus(U_0\cup U_1)$, so $|M|\ge5$.
For each $b$ with $t_b>0$, construct nested sets exactly as above, except that
the first backward expansion has size $4$ and each earlier round adds one vertex.
At a first-expansion round, choose the external connector in $M$, distinct from
any other connector in that snapshot; at earlier rounds use the fresh pool vertex
that is added to the nested cone.

Use the cubic protected-block lemma for each active target.  A connector has
two incident block edges when the protected-set size is even and one incident
block edge when that size is odd, so it requires residual degree $1$ or $2$
inside the complement.  All non-connectors require degree $3$ there.
At every round the complement has at least five vertices.  It remains only
to check the parity condition in the cubic complement-completion lemma.
Let $S$ be the union of the protected blocks and let $e_{\rm cut}$ be the
number of already fixed block-to-connector edges.  Since every vertex of $S$
already has degree $3$,
\[
3|S|=2|E(S)|+e_{\rm cut},
\]
so
\[
e_{\rm cut}\equiv |S|\pmod 2.
\]
The required degree sum inside the complement $R=V\setminus S$ is
\[
3|R|-e_{\rm cut}
\equiv |R|-|S|
\equiv n
\equiv0\pmod2,
\]
because $n$ is even.  Hence the cubic complement-completion lemma applies.

As in the quartic case, induction on the backward-cone recursion gives the
prescribed nested sets exactly for positive horizons, while a zero-horizon cone is
the singleton target.  The two time-zero cones are therefore disjoint.
\end{proof}

\section{Near-exact known-size lower bounds}

\begin{theorem}[Quartic known-size lower bound]
For every $n\ge15$ and $P\ge1$,
\[
\boxed{
S_n(n,P)\ge
2+\left\lceil\frac{n-13}{P}\right\rceil.
}
\]
\end{theorem}

\begin{proof}
Fix a correct algorithm with state set $Q$, and let $\tau_0,\tau_1$ be the two
homogeneous first-final times on $4$-regular schedules.

If $\tau_0+\tau_1\le n-14$, the quartic slow-cone proposition gives a connected
$4$-regular schedule with two targets whose time-zero backward cones are disjoint.
Assign input $0$ to every node in the first cone and input $1$ to every node in
the second cone; assign the remaining inputs arbitrarily.  By dynamic causal-cone
locality, target $b$ reproduces its homogeneous input-$b$ history through time
$\tau_b$ and therefore reaches an immutable final output $b$.  The two opposite
final decisions coexist, contradicting agreement.  Hence
\[
\tau_0+\tau_1\ge n-13.
\]
Regular two-orbit packing now gives
\[
P|Q|\ge n-13+2P,
\]
which is equivalent to the stated bound.
\end{proof}

\begin{theorem}[Even-size cubic refinement]
For every even $n\ge12$ and $P\ge1$,
\[
\boxed{
S_n(n,P)\ge
2+\left\lceil\frac{n-10}{P}\right\rceil.
}
\]
\end{theorem}

\begin{proof}
Use the cubic slow-cone proposition.  Correctness implies
\[
\tau_0+\tau_1\ge n-10
\]
for the two homogeneous first-final times on cubic schedules.  Applying regular
two-orbit packing yields
\[
P|Q|\ge n-10+2P,
\]
and hence the result.
\end{proof}

Combining the quartic lower bound with the existing block-counter upper bound gives,
for $n\ge15$,
\[
2+\left\lceil\frac{n-13}{P}\right\rceil
\le S_n(n,P)\le
2+\left\lceil\frac{n-1}{P}\right\rceil.
\]
For even $n\ge12$, the left side improves by replacing $n-13$ with $n-10$.

\section{Exact three-state phase threshold}

\begin{lemma}[Three-state normal form]
Suppose a correct algorithm has exactly three persistent states.  Then, after
possibly exchanging the names $0$ and $1$, its states may be written
\[
Q=\{\Fzero,\Live,\Fone\},
\]
where $\Fzero,\Fone$ are distinct immutable final states with outputs $0,1$,
$\Live$ is the unique nonfinal state, and
\[
\iota(0)=\Fzero,\qquad \iota(1)=\Live.
\]
\end{lemma}

\begin{proof}
Unanimity requires that some reachable final state output $0$ and some distinct
reachable final state output $1$.  Thus at least two states are final.
A correct mixed execution cannot start with both input values already in opposite
immutable final states, so a nonfinal state is also necessary.  With only three
states, there are therefore exactly two relevant final states and exactly one
nonfinal state $\Live$.

An input-$b$ node cannot initialize in the opposite final state, because that would
already violate unanimity in the homogeneous input-$b$ execution.  If both input
values initialized in $\Live$, their homogeneous executions would be identical and
could not terminate with opposite outputs.  If both initialized in their respective
final states, a mixed input would start with opposite immutable decisions.
Hence exactly one input initializes final and the other initializes in $\Live$.
Exchange the bit names if necessary.
\end{proof}

\begin{lemma}[Path symmetry before the first one-finalization]
Run a three-state algorithm in the all-one execution on the static path
\[
v_0-v_1-\cdots-v_{n-1}.
\]
Before the first transition into $\Fone$, every node is in state $\Live$.
At every such time, the two endpoints have identical local observations and all
interior vertices have identical local observations.  Consequently the first
$\Fone$ transition, when it occurs, is simultaneous on the entire endpoint class,
the entire interior class, or both.
\end{lemma}

\begin{proof}
No node may enter $\Fzero$ in the all-one execution, because $\Fzero$ is immutable
with output $0$.  Since $\Live$ is the only nonfinal state, every node therefore
remains in $\Live$ until the first $\Fone$ transition.  The two endpoints each see
one $\Live$ neighbor, while every interior node sees two $\Live$ neighbors.
The phase is common and the transition rule is anonymous and deterministic.
\end{proof}


\begin{lemma}[Forced maximal zero wave at the boundary]
Assume the three-state normal form and suppose that, in the all-one execution
on $P_n$, the first final-one event is interior-only at time
\[
P=n-2.
\]
In the mixed execution with $v_0$ initially zero and all other inputs one,
correctness forces, for every $0\le j\le P-2$,
\[
q_{v_j}(j)=\Fzero,
\qquad
q_{v_{j+1}}(j)=\Live,
\]
and the propagation-round-$j$ transition at the frontier satisfies
\[
\delta\!\left(
\Live,\phi_j,
\text{one broadcast from }\Fzero
\text{ and one broadcast from }\Live
\right)=\Fzero.
\]
Consequently
\[
\sigma(\Fzero,\phi_j)\ne\sigma(\Live,\phi_j)
\qquad(0\le j\le P-2).
\]
Moreover the same inequality also holds at phase $P-1$.
\end{lemma}

\begin{proof}
The zero source $v_0$ is permanently in $\Fzero$, and correctness forbids any
node from entering $\Fone$ in the same mixed execution.  The only remaining
state is therefore $\Live$.

In the all-one execution, $v_{n-2}$ enters $\Fone$ at time $P=n-2$.
For the mixed execution to avoid that same transition, the causal difference
originating at $v_0$ must affect the round-$P-1$ observation of $v_{n-2}$.
The distance from $v_0$ to $v_{n-3}$ is exactly
\[
n-3=P-1.
\]
One-hop-per-round causality leaves no slack: a delay in any one of rounds
$0,\ldots,P-2$ cannot be recovered later.  Hence at time $j$ the influence
frontier must already have reached $v_j$, and because $\Fone$ is forbidden,
the differing state there must be $\Fzero$.  The next vertex $v_{j+1}$ is
still outside the causal influence and remains in $\Live$ at the beginning
of round $j$.  To advance the frontier, its next state must be $\Fzero$.
This proves the displayed transition.

If the two broadcasts at phase $\phi_j$ were equal, the frontier vertex
$v_{j+1}$ would receive exactly the same count vector as its all-one
counterpart and could not make the required different transition.  Thus the
broadcasts differ for $0\le j\le P-2$.

At the last phase $P-1$, the vertex $v_{n-2}$ begins the round in $\Live$.
Its neighbors are $v_{n-3}\in\Fzero$ and $v_{n-1}\in\Live$.  If the two
broadcasts were equal at that phase, its received count vector would again
equal the all-one count vector, forcing the forbidden transition into
$\Fone$.  Hence the inequality also holds at phase $P-1$.
\end{proof}

\begin{theorem}[Exact three-state threshold]
For every $n\ge4$ and $P\ge1$,
\[
\boxed{
S_n(n,P)=3
\quad\Longleftrightarrow\quad
P\ge n-1.
}
\]
\end{theorem}

\begin{proof}
Any correct algorithm needs at least three persistent states: unanimity requires
distinct immutable final states for outputs $0$ and $1$, while mixed-input
correctness requires a nonfinal state.  The existing block-counter construction
with horizon $H=n-1$ uses one live state and two final states when $P\ge n-1$.
Hence $S_n(n,P)=3$ in that regime.

For the converse, suppose that a correct three-state algorithm exists with
$P\le n-2$.  Put it in the normal form
$\{\Fzero,\Live,\Fone\}$.

Consider the all-one execution on the static path $P_n$, and let $T$ be the first
time at which some node enters $\Fone$.  If no $\Fone$ appears during the first
$P$ propagation rounds, then at time $P$ every node is again in $\Live$ and the
phase is again $0$.  The entire global configuration has repeated, so the execution
would never terminate.  Hence
\[
T\le P\le n-2.
\]

If an endpoint is among the first nodes entering $\Fone$, put a single zero at the
opposite endpoint.  Its distance is $n-1>T$, so the first endpoint has exactly the
same local history through time $T$ as in the all-one execution and enters
$\Fone$, while the zero endpoint is already in the immutable state $\Fzero$.
This violates agreement.

Thus the first event is interior-only.  By path symmetry all interior vertices
enter $\Fone$ simultaneously.  If $T<n-2$, put a single zero at $v_0$.
The vertex $v_{n-2}$ is at distance $n-2>T$ from that zero, so its local history
through time $T$ is unchanged and it enters $\Fone$, again conflicting with the
immutable $\Fzero$ at $v_0$.  Therefore the only remaining case is
\[
T=P=n-2,
\]
with all interior vertices entering $\Fone$ at time $P$ and both endpoints still
in $\Live$.

For each phase $\phi$, define the broadcast bits
\[
a_\phi=\sigma(\Fzero,\phi),\qquad
\ell_\phi=\sigma(\Live,\phi),\qquad
b_\phi=\sigma(\Fone,\phi).
\]

Apply the forced maximal zero-wave lemma.  In terms of the broadcast
notation just introduced, it gives
\[
a_\phi\ne\ell_\phi
\qquad\text{for every }\phi\in\{0,\ldots,P-1\},
\]
and, for phases $0,\ldots,P-2$, it also gives the concrete frontier transition
$\Live\to\Fzero$ when an interior live node sees one $\Fzero$ neighbor and one
$\Live$ neighbor.

Return to the all-one execution.  At time $P$ every interior vertex is in the
immutable state $\Fone$, while the two endpoints are still in $\Live$.
Termination requires each endpoint eventually to enter $\Fone$.
Let $r\ge P$ be a propagation round at whose end an endpoint first enters
$\Fone$, and put
\[
\psi=r\bmod P.
\]
At the beginning of round $r$ that endpoint is in $\Live$ and its unique neighbor
is in $\Fone$.  During the first phase-$\psi$ round of the execution, the same
endpoint was in $\Live$ with a unique $\Live$ neighbor and remained in $\Live$.
Therefore
\[
b_\psi\ne\ell_\psi.
\]
Since the broadcast alphabet is binary and also
$a_\psi\ne\ell_\psi$, necessarily
\[
a_\psi=b_\psi.
\]

Finally, take a static path and choose a target endpoint $u$.
Place a single zero at the vertex whose distance from $u$ is $\psi+1$; all other
inputs are one.  The forced frontier transitions above propagate $\Fzero$ toward
$u$ at one edge per round during phases $0,\ldots,\psi-1$.  Hence at the beginning
of propagation round $\psi$, the unique neighbor of $u$ is in $\Fzero$, while
$u$ itself is still in $\Live$.  The endpoint therefore sees the same received
bit $a_\psi=b_\psi$ that caused an all-one endpoint with an $\Fone$ neighbor to
enter $\Fone$ at phase $\psi$.  Determinism forces $u$ into $\Fone$.
The source zero remains permanently in $\Fzero$, contradicting agreement.

Hence no three-state algorithm exists when $P\le n-2$.
\end{proof}


\section{Construction audit}

The graph lemmas above are explicit and constitute the proof.  As a separate implementation
audit, the constructions were independently instantiated over finite parameter ranges and
their claimed invariants were recomputed from the generated snapshots.

For quartic schedules, every ordered pair $(t_0,t_1)$ with
\[
15\le n\le40,\qquad t_0,t_1\ge0,\qquad t_0+t_1\le n-14
\]
was checked, for a total of $3653$ cases.  For cubic schedules, every ordered pair with
\[
12\le n\le40,\quad n\text{ even},\qquad
t_0,t_1\ge0,\qquad t_0+t_1\le n-11
\]
was checked, for a total of $2600$ cases.  Across all $6253$ complete schedules, every
snapshot was simple, connected, and exactly regular of the intended degree; recomputed
backward causal cones matched the prescribed nested cones; and the two time-zero cones were
disjoint.  No construction failure was observed.

These checks are not part of the mathematical proof and do not establish the theorem for
unbounded $n$.  Their purpose is to test the indexing and explicit graph constructions used
in the proof.

\section{Consequences and remaining gap}

Whenever $n/P\to\infty$ (in particular, for fixed or sufficiently slowly growing $P$), the quartic theorem and the block-counter upper bound give
\[
S_n(n,P)=\frac{n}{P}+O(1)
\]
with leading constant one.  At $P=1$ the general bounds become
\[
n-11\le S_n(n,1)\le n+1
\qquad(n\ge15),
\]
while even $n\ge12$ improves the lower side to
\[
n-8\le S_n(n,1)\le n+1.
\]

The exact three-state result identifies the first saturation boundary:
\[
P_{\min}^{(3)}(n)=n-1.
\]
It does not by itself imply the full feasible-region formula for larger state budgets.
In particular, we leave open whether
\[
S_n(n,P)=2+\left\lceil\frac{n-1}{P}\right\rceil
\]
holds for all $n$ and $P$, or whether dynamic topology can be used by algorithms as a
distributed spatial memory to beat the block-counter upper construction for intermediate
state budgets.

\section{Claim boundary and limitations}

The paper does not claim novelty for finite-state pumping, deterministic trajectory
packing in general, causal-cone locality, regular graph constructions, monotone flooding,
periodic clocks, or generic time--space tradeoffs.  The claimed contribution is the
task-specific use of regular dynamic witnesses to sharpen the known-size state lower bound
to within an additive constant of the existing upper numerator, together with the exact
three-state phase threshold.

The model assumes deterministic anonymous synchronous leaderless
$1$-interval-connected dynamic networks; a free globally aligned period-$P$ phase; output
and finality encoded in persistent state rather than in the phase; final nodes remaining
present in communication; and one-bit broadcast-counting communication.  Randomization,
identifiers, ports, leaders, Byzantine faults, asynchronous execution, and
disappearance-on-finalization semantics are outside scope.

\section{Conclusion}

Regularity lets a dynamic adversary preserve the homogeneous state orbit while changing
the geometry of causal influence.  This separates two roles that were entangled in the
previous static-cycle lower bound: symmetry of the unanimous executions and propagation
distance in a mixed execution.  Explicit quartic and cubic schedules then force the sum of
the two homogeneous first-final times to be $n-O(1)$, and two-orbit phase packing converts
that waiting into a near-exact persistent-state lower bound.

At the minimum state budget, the structure is sharper.  Three states leave only one live
state, and on a static path this forces a maximal-speed zero wave at the critical boundary.
The one-bit alphabet then makes final-zero and final-one communication indistinguishable at
some phase, yielding the exact threshold $P=n-1$.

The remaining problem is no longer a factor-two uncertainty.  It is whether the final
additive gap can be closed for arbitrary state budgets.

\begin{thebibliography}{99}

\bibitem{yonezu2026statephase}
R. Yonezu,
``Internal Time as Local Memory: State--Phase Tradeoffs for Explicit Termination in Anonymous Dynamic Binary Consensus,''
Zenodo preprint, version 2.0.0, September 2, 2026.
DOI: \href{https://doi.org/10.5281/zenodo.22228873}{10.5281/zenodo.22228873}.

\bibitem{yonezu2026twoorbit}
R. Yonezu,
``Two-Orbit Packing for Exact State Complexity of Explicit Binary Consensus in Anonymous Dynamic Networks with Periodic Time,''
Zenodo preprint, version 1.0.0, September 6, 2026.
DOI: \href{https://doi.org/10.5281/zenodo.22538052}{10.5281/zenodo.22538052}.
This preprint is the direct predecessor of the present work and is not treated here as
independent prior art for novelty purposes.

\bibitem{blanc2026}
T. Blanc, G. A. Di Luna, and G. Viglietta,
``Computing in Anonymous Dynamic Networks with One-Bit Communications,''
arXiv:2607.08358, 2026.

\bibitem{parzych2026}
G. Parzych and J. J. Daymude,
``Memory Lower Bounds and Impossibility Results for Anonymous Dynamic Broadcast,''
\emph{Distributed Computing}, vol. 39, no. 3, Art. 18, 2026.
DOI: \href{https://doi.org/10.1007/s00446-026-00511-4}{10.1007/s00446-026-00511-4}.

\bibitem{diluna2024}
G. A. Di Luna and G. Viglietta,
``Universal Finite-State and Self-Stabilizing Computation in Anonymous Dynamic Networks,''
in \emph{OPODIS 2024}, LIPIcs 324, Art. 10, 2024.
DOI: \href{https://doi.org/10.4230/LIPIcs.OPODIS.2024.10}{10.4230/LIPIcs.OPODIS.2024.10}.

\bibitem{turau2026}
V. Turau,
``Broadcasts in Anonymous, Dynamic Networks: A New Algorithm and Impossibility Results,''
in \emph{SAND 2026}, LIPIcs 373, Art. 6, 2026.
DOI: \href{https://doi.org/10.4230/LIPIcs.SAND.2026.6}{10.4230/LIPIcs.SAND.2026.6}.

\bibitem{monterno2023}
L. Penet de Monterno, B. Charron-Bost, and S. Merz,
``Synchronization modulo $P$ in Dynamic Networks,''
\emph{Theoretical Computer Science}, vol. 942, pp. 200--212, 2023.

\bibitem{charron2022}
B. Charron-Bost and L. Penet de Monterno,
``Self-Stabilizing Clock Synchronization in Dynamic Networks,''
in \emph{OPODIS 2022}, LIPIcs 253, Art. 28, 2022.
DOI: \href{https://doi.org/10.4230/LIPIcs.OPODIS.2022.28}{10.4230/LIPIcs.OPODIS.2022.28}.

\end{thebibliography}

\end{document}
